\section{A balanced minimal forced-family feed has no far-in source}
\label{sec:balanced-forced-feed-minimal-refinement}

Here $D$ is finite, nonempty, strongly connected, all-deficient and
minimal under deletion of arbitrary nonempty arc sets on its fixed
vertex set. Let $T_*$ be a maximum protection-forced reference completion.
Fix any global median feed $f$ satisfying the \emph{balanced} condition
$t_D(f)=g_D(f)$. Put $p=|P(f)|$, $\mu=\mu(f)$,
$g=g_D(f)$, $Q=N_D^{--}(f)$, and $C=\partial P(f)$.

\begin{proposition}[The balanced-feed empty far-in refinement]
\label{prop:balanced-forced-feed-three-branches}
\label{prop:balanced-forced-feed-far-in-empty}
One has $\mu\ge1$ and $P(f)=\varnothing$. The only remaining
numerical shapes are
\[
 (g,t_D(f),p,\eta(f))\in
 \bigl\{(1,1,0,2\mu),\ (2,2,0,2\mu+1)\bigr\}.
\]
Thus every other original vertex reaches $f$ within at most two steps.
This is an incoming-reachability statement, not an assertion that $f$
has the original second-neighborhood property or is an actual BAD source.
\end{proposition}
\begin{proof}
The general balanced-feed theorem
\ref{thm:forced-balanced-feed-boundary} gives
$\mu\ge1$ and $\mathcal M(f)\subseteq Q\setminus C$.
The strong deletion consequence \ref{lem:gap} gives $g\in\{1,2\}$.
The residual identity is
\[
 p=2\mu+g-1-\eta(f).
\]
If $P(f)=\varnothing$, use only this identity: it gives the first
branch when $g=1$ and the second when $g=2$. The nonempty-$P$
packet estimates below are \emph{not} applied in this case.

Assume $P(f)\ne\varnothing$. In the established partner-packet notation
let $J=Q\setminus C$, $B=J\cap\mathcal M(f)$,
$k=|B|$, and $s=|J\cap N_D^+(f)|$. The containment above implies
$B=\mathcal M(f)$ and $k=\mu$; the target-bundle partition gives
$|J|=k+s=\mu+s$. Lemma~\ref{lem:partner-head-budget}, whose stated
hypothesis includes $P(f)\ne\varnothing$, gives a genuine set $L$ with
\[
 0\le|L|\le
 \begin{cases}
   \eta(f)-2\mu,&s=0,\\
   \eta(f)-g-2\mu,&s>0.
 \end{cases}
\]
Equivalently these two lower bounds on $\eta$ follow directly from
the original target-bundle refinements
$\eta\ge2|J|$ for $s=0$, and
$\eta\ge g+2(|J|-s)$ for $s>0$.

If $s>0$, substituting $\eta\ge g+2\mu$ into the residual identity
gives $p\le-1$, impossible. Hence $s=0$ and $\eta\ge2\mu$,
so $1\le p\le g-1$. Since $g\in\{1,2\}$, this forces
$g=2$, $p=1$, and then $\eta=2\mu$. Also $|J|=\mu$ and
$\mathcal M(f)\subseteq J$, so $J=\mathcal M(f)$.

Let $P(f)=\{z\}$. A member of $P(f)$ cannot point to $f$,
and $P(f)$ is disjoint from $Q$. As every missing partner is in $Q$,
the vertex $z$ is not a missing partner either. Therefore the original
direction is $f\to z$. Since $P(f)$ is a singleton,
$C=N_D^+(z)$. No member of $C$ can point to $f$, for that would
give $z$ a two-step path to $f$. Nor can a member of $C$ be missing
from $f$, since $\mathcal M(f)=Q\setminus C$ is disjoint from $C$.
Consequently $C\subseteq N_D^+(f)$. The target-bundle partition
\[
 V(D)=\{f\}\,\dot\cup\,N_D^-(f)\,\dot\cup\,P(f)
       \,\dot\cup\,C\,\dot\cup\,(Q\setminus C)
\]
and $Q\setminus C=\mathcal M(f)$ exclude all other first neighbors,
so
\[
 N_D^+(f)=\{z\}\,\dot\cup\,C
          =\{z\}\,\dot\cup\,N_D^+(z).
\]
Each $z\to c$ with $c\in C$ is original, so no $c\in C$ points
to $z$. Delete just the original arc $f\to z$, and call the result
$D'$. The removed intermediate $z$ had only heads in $C$, still
original first neighbors of $f$, so removing it destroys no old exact
second target of $f$. The removed first target $z$ cannot become
exact second, since no retained first intermediate $c\in C$ points
to $z$. Thus
\[
 N_{D'}^{++}(f)=N_D^{++}(f),\qquad
 g_{D'}(f)=g_D(f)-1=1.
\]
Every other vertex has unchanged first neighborhood and can only lose
second targets. Hence $D'$ is still all-deficient, contradicting
arc-set deletion minimality. This excludes the only possible
nonempty-$P$ branch, so $P(f)=\varnothing$ and the two stated formulas
follow.
\end{proof}

\begin{corollary}[Nonempty far-in forces strict head surplus]
\label{cor:forced-minimal-far-in-head-surplus}
In the same minimal-counterexample class, every global median feed $f$
of $T_*$ with $P(f)\ne\varnothing$ satisfies
\[
 t_D(f)\ge g_D(f)+1.
\]
\end{corollary}
\begin{proof}
The general feed inequality gives $t_D(f)\ge g_D(f)>0$.
Equality would imply $P(f)=\varnothing$ by the proposition.
Integrality gives the strict-surplus bound.
\end{proof}

\paragraph{What remains unresolved.}
The proposition does not exclude either of its two remaining branches.
In the $g=1$ balanced branch, $\eta(f)=2\mu(f)\ge2$, but there may still be
arbitrarily many forced incoming missing partners and reversed unique-GOOD
second arcs. In the $g=2$ balanced branch, $\eta(f)=2\mu(f)+1\ge3$;
this branch is also unresolved. Feeds with strict head surplus
$t_D(f)>g_D(f)$ are outside this refinement, and their far-in sets
need not be small. In particular this is not a universal claim
$P(f)=\varnothing$ for every
maximal forced-family feed, nor a reduction of the full SNC to finitely
many graph orders or residual allocations.
